\daytitle{7}{连接体、斜面与摩擦综合}

\begin{learningbox}{本节导读}
本节是运动学和力学的第一次综合实战。连接体题先判断多个物体之间的约束，再决定用整体法求共同运动，还是用隔离法求内部作用力。斜面临界和传送带模型都要回到同一条主线：相对运动趋势决定摩擦方向，约束关系决定加速度关系。
\end{learningbox}

\textbf{本节读题路线：} 先判断多个物体是否有共同加速度，以及绳长、接触或相对运动给出什么约束；只求整体加速度先用整体法，要求张力或支持力再隔离。斜面临界题先看摩擦趋势，传送带题要按“未共速、共速后能否维持共速”分阶段。

\section*{一、学习目标}

\begin{enumerate}[itemsep=0.65em,topsep=0.4em,parsep=0pt]
  \item 能根据所求量选择整体法或隔离法，并写出对应的动力学方程。
  \item 会用整体法求连接体的加速度，用隔离法求内力（绳中张力）。
  \item 能推导斜面摩擦临界条件 $\tan\theta = \mu_s$，并说明它的应用边界。
  \item 能按“未共速—共速瞬间—共速后”分析水平和倾斜传送带，并检查传送带长度是否足够。
  \item 能从整根轻绳长度不变出发，建立不同运动方向之间的速度、加速度约束。
  \item 能综合运用受力分析、约束关系和牛顿定律解决多物体问题。
\end{enumerate}

\section*{二、初高中衔接}

初中物理中大部分问题只涉及单个物体。高中物理中经常出现两个或更多物体连接在一起运动的情况——这就是\textbf{连接体问题}。

连接体问题的核心思路只有两条：
\begin{enumerate}
  \item \textbf{整体法：} 把几个物体看成一个整体，只分析整体受到的外力，求出共同加速度。
  \item \textbf{隔离法：} 把其中一个物体分离出来单独分析，求内部相互作用力（如绳的张力）。
\end{enumerate}

这两个方法交替使用，就是解决连接体问题的标准流程。

\textbf{科学发展视角：} 工厂传送带、电梯曳引系统、起重机吊装和机器人多关节联动，都不是单个物体孤立运动，而是多个部件通过绳、带、接触面或关节形成约束。整体法用于观察系统的外部表现，隔离法用于寻找内部关键受力，两者结合才符合真实工程分析。

\section*{三、核心概念}

\begin{conceptbox}{整体法与隔离法}

\textbf{整体法：}
\begin{itemize}
  \item 适用范围：多个物体具有相同的加速度矢量（相对静止并一起平动）。
  \item 操作：把多个物体看作一个整体，只画整体受到的外力（不画内部物体之间的力）。
  \item 列方程：$F_{\text{合外}} = (m_1 + m_2 + \cdots) a$
\end{itemize}

\textbf{隔离法：}
\begin{itemize}
  \item 适用范围：需要求系统内部物体间的相互作用力（如绳张力、物体间支持力）。
  \item 操作：从系统中分离出一个物体单独做受力分析，列出牛顿第二定律方程。
  \item 列方程：$F_{i,\text{合}} = m_i a$（共同运动时，$a$ 与整体法中的 $a$ 相同）
\end{itemize}

\textbf{选择策略：}
\begin{itemize}
  \item 只求加速度 $\to$ 整体法。
  \item 求内力 $\to$ 先整体法求加速度，再隔离法求内力。
  \item 各物体加速度方向或大小不同 $\to$ 不能直接把所有质量相加套一个方向的整体法，通常要逐个隔离。
\end{itemize}

\textbf{约束关系从哪里来：} 轻绳不伸长，所以绳长是常量。写出绳长关系，才能判断两个物体的速度、加速度大小是否相等。若绳子斜拉物体，绳长的变化只由物体速度沿绳方向的分量决定，不能直接把物体速度等同于绳端速度。

\pagebreak
\begin{center}
\begin{tikzpicture}[scale=0.95,>=stealth]
  \coordinate (O) at (1.75,0);
  \def\r{0.32}

  % 墙角：滑轮圆心固定在桌边墙角处
  \draw[thick] (-2.0,0) -- (O) -- (1.75,-2.35);
  \foreach \x in {-1.8,-1.45,...,1.45} {
    \draw[gray!60] (\x,0) -- ++(0.18,-0.12);
  }
  \foreach \y in {-0.25,-0.60,...,-2.15} {
    \draw[gray!60] (1.75,\y) -- ++(0.12,0.18);
  }

  \draw[fill=white,thick] (-1.5,0) rectangle (-0.6,0.55);
  \node at (-1.05,0.28) {$A$};

  \draw[fill=white,thick] (1.87,-1.85) rectangle (2.45,-1.0);
  \node at (2.16,-1.43) {$B$};

  \draw[thick,fill=white] (O) circle (\r);
  \fill (O) circle (0.025);

  % 一根连续轻绳：A 端--滑轮--B 端
  \draw[thick] (-0.6,\r) -- ($(O)+(0,\r)$)
    arc[start angle=90,end angle=0,radius=\r]
    -- ($(O)+(\r,-1.0)$);

  \draw[<->,blue] (-0.6,0.88) -- ($(O)+(0,0.88)$) node[midway,above] {$l_1$};
  \draw[<->,red] ($(O)+(0.62,0)$) -- ($(O)+(0.62,-1.0)$) node[midway,right] {$l_2$};
  \draw[->,thick,blue] (-1.05,-0.18) -- (-0.25,-0.18) node[right] {$a_A$};
  \draw[->,thick,red] (2.16,-1.85) -- (2.16,-2.48) node[midway,right] {$a_B$};

  \node[align=center] at (-0.25,-3.35)
    {$l_1+l_2+\text{定长部分}=L$\\[2pt]
     $\Delta l_1+\Delta l_2=0\Rightarrow |a_A|=|a_B|$};
\end{tikzpicture}


\figcaption{轻绳不伸长：由绳长不变推出连接体的运动约束}\label{fig:rope-constraint}
\end{center}

\end{conceptbox}

\begin{conceptbox}{摩擦临界问题}

\textbf{临界条件：} 物体处于将要滑动的极限静止状态时，静摩擦力达到最大值。

\textbf{斜面摩擦临界：} 对只受重力、支持力和静摩擦力的物块，若没有其他沿斜面方向的外力，刚好不下滑的条件为：
\[
mg\sin\theta = f_{\max} = \mu_s mg\cos\theta
\]
\[
\boxed{\tan\theta = \mu_s}
\]

\begin{itemize}
  \item $\theta < \arctan\mu_s$：若物块初始静止，则可保持静止（$f_s = mg\sin\theta < f_{\max}$）。
  \item $\theta = \arctan\mu_s$：物块可处于极限静止状态，此时 $f_s=f_{\max}$；倾角稍增或受到沿斜面向下的微小扰动后将滑动。
  \item $\theta > \arctan\mu_s$：若由静止释放，物块将下滑（滑动摩擦力 $f_k = \mu_k mg\cos\theta$）。
\end{itemize}

\begin{center}
% TikZ 图：斜面摩擦临界
% Day 07：连接体、斜面与摩擦综合
\begin{tikzpicture}[scale=0.92,>=Latex]
  % 斜面：右下方为可能下滑方向。
  \draw[thick] (0,0) -- (5,0) -- (0,2) -- cycle;
  \draw[thick] (0,2) -- (5,0);
  \draw (4.25,0) arc (180:158.2:0.75);
  \node at (4.1,0.25) {$\theta$};

  \coordinate (P) at (2.35,1.06);
  \coordinate (C) at ($(P)+(68.2:0.42)$);

  \begin{scope}[shift={(P)},rotate=-21.8]
    \draw[thick] (-0.6,0) rectangle (0.6,0.85);
    \node at (0,0.43) {$m$};
  \end{scope}

  % 真实力
  \draw[->,thick,blue] (C) -- ++(0,-1.25) node[below] {$G$};
  \draw[->,thick,blue] (C) -- ++(68.2:1.05) node[above] {$N$};
  \draw[->,thick,red] (C) -- ++(158.2:1.15) node[above] {$f_s=f_{\max}$};

  % 分力
  \draw[->,thick,red,dashed] (C) -- ++(-21.8:1.25) node[right] {$G\sin\theta$};
  \draw[->,thick,blue,dashed] (C) -- ++(-111.8:0.85) node[left] {$G\cos\theta$};

  \node[font=\footnotesize] at (2.65,-0.45) {临界：$G\sin\theta=\mu_s N$};
\end{tikzpicture}


\figcaption{斜面摩擦临界：$G\sin\theta = \mu_s N$ 时物块处于极限静止状态}\label{fig:critical-incline}
\end{center}

\end{conceptbox}

\section*{四、公式与适用条件}

\textbf{方法使用前检查：}
\begin{enumerate}
  \item 多个物体是否有共同加速度？若没有，不能直接把质量相加。
  \item 只求整体加速度时用整体法；要求张力、支持力等内力时再隔离。
  \item 轻绳、轻杆、接触面和传送带给出的约束关系要先写清。
  \item 摩擦临界和传送带题都要先判断相对运动趋势，再决定摩擦方向。
\end{enumerate}

\begin{learningbox}{连接体的系统边界分析清单}
整体法的本质是先画一个“系统边界”。边界里面的物体合成一个整体，边界内部彼此之间的作用力会成对抵消；边界外的物体对系统的作用力才进入整体方程。

\begin{center}
\begin{tabular}{
  >{\raggedright\arraybackslash}p{0.22\textwidth}
  >{\raggedright\arraybackslash}p{0.34\textwidth}
  >{\raggedright\arraybackslash}p{0.30\textwidth}}
\toprule
先问什么 & 判断结果 & 对应做法 \\
\midrule
把谁圈成系统？ & $A+B$、人+船、物块+小车等 & 写总质量 $m_{\text{总}}$ \\
哪些力来自系统外？ & 拉力、重力、地面摩擦、墙面支持等 & 写进整体方程 \\
哪些力在系统内部？ & 绳张力、两物体间支持力等 & 整体法中不写 \\
是否要求内部力？ & 要求张力、压力、支持力 & 再隔离其中一个物体 \\
\bottomrule
\end{tabular}
\end{center}

\textbf{判断标准：} 一个力是否写进整体方程，不看它“有没有出现”，而看施力物体是否在系统边界外。改变系统边界后，同一个力可能从内力变成外力，所以整体法和隔离法必须说清研究对象。
\end{learningbox}

\begin{center}
\begin{tabular}{lll}
\toprule
方法 & 适用场景 & 核心关系 \\
\midrule
整体法 & 共同加速度系统 & $F_{\text{合外}} = m_{\text{总}} a$ \\
隔离法 & 求内部相互作用力 & $F_{i,\text{合}} = m_i a$ \\
斜面临界 & 无沿斜面外力 & $\tan\theta = \mu_s$ \\
水平传送带 & 判断摩擦方向 & 共速且无水平外力时 $f=0$ \\
\bottomrule
\end{tabular}
\end{center}

\section*{五、典型模型}

\subsection*{模型 1：水平面连接体}

两个物体用轻绳连接在粗糙水平面上，前方物体 $A$ 受水平拉力 $F$，后方物体 $B$ 被绳拉着一起向右运动。

\begin{center}
% TikZ 图：水平连接体整体法
% Day 07：连接体、斜面与摩擦综合
\begin{tikzpicture}[scale=0.9,>=Latex]
  % 地面
  \draw[thick] (-1.0,0) -- (6.3,0);

  % 后方 B，前方 A
  \draw[thick] (0,0) rectangle (1.5,1.1);
  \node at (0.75,0.55) {$m_B$};
  \draw[thick] (3.0,0) rectangle (4.5,1.1);
  \node at (3.75,0.55) {$m_A$};

  % 轻绳
  \draw[thick] (1.5,0.55) -- (3.0,0.55);
  \node[above] at (2.25,0.55) {轻绳};

  % 外力和外摩擦
  \draw[->,thick,red] (4.5,0.55) -- (5.8,0.55) node[above] {$F$};
  \draw[->,thick,red] (0,0.25) -- (-0.9,0.25) node[above] {$f_B$};
  \draw[->,thick,red] (3.0,0.18) -- (2.1,0.18) node[below] {$f_A$};
  \draw[->,thick,blue] (2.6,1.55) -- (3.7,1.55) node[right] {$a$};

  % 整体法示意
  \draw[thick] (0,-0.3) -- (4.5,-0.3);
  \draw[thick] (0,-0.2) -- (0,-0.4);
  \draw[thick] (4.5,-0.2) -- (4.5,-0.4);
  \node[below] at (2.25,-0.62) {整体：$F-(f_A+f_B)=(m_A+m_B)a$};
\end{tikzpicture}


\figcaption{连接体整体法：整体 $F - f_{\text{总}} = (m_A + m_B) a$}\label{fig:connected-body}
\end{center}

\noindent 图中只保留水平方向上参与列方程的力；重力和支持力在竖直方向平衡，整体法中不写绳的张力这种内力。

\noindent 隔离 $B$ 求张力：
\begin{center}
% TikZ 图：连接体隔离法（分析B受力）
% Day 07：连接体、斜面与摩擦综合
\begin{tikzpicture}[scale=0.9,>=Latex]
  \draw[thick] (-1,0) -- (4,0);
  % B 物块
  \draw[thick] (0,0) rectangle (1.5,1.2);
  \node at (0.75,0.6) {$m_B$};
  % 绳拉力 T
  \draw[->,thick,blue] (0,0.6) -- (-1,0.6) node[above] {$T$};
  % 摩擦力
  \draw[->,thick,red] (1.5,0.3) -- (2.5,0.3) node[above] {$f_B$};
  % 隔离法标注
  \draw[decorate,decoration={brace,amplitude=4pt}] (0,-0.3) -- (1.5,-0.3);
  \node[below] at (0.75,-0.6) {隔离 $m_B$：$T - f_B = m_B a$};
\end{tikzpicture}


\figcaption{隔离 $B$：$T - f_B = m_B a$}\label{fig:isolated-body}
\end{center}

\subsection*{模型 2：斜面摩擦临界}

这个模型必须先限定为：物块初始静止，只受重力、支持力和摩擦力，且没有其他沿斜面方向的外力。判断步骤为：
\begin{enumerate}
  \item 假设物块有沿斜面向下的运动趋势，静摩擦力沿斜面向上。
  \item 先由 $N=mg\cos\theta$ 求最大静摩擦力 $f_{\max}=\mu_sN$。
  \item 比较 $mg\sin\theta$ 与 $f_{\max}$：前者较小时可静止，相等时为极限静止，前者较大时物块下滑。
\end{enumerate}
只有在上述条件下，临界式才能化为 $\tan\theta=\mu_s$；有拉力、推力或其他约束时，应重新列沿斜面方向的平衡方程。

\subsection*{模型 3：水平传送带}

物块轻放在以 $v_{\text{带}}$ 匀速运动的水平传送带上。刚放上时
$v_{\text{物}}=0$，物块相对带面向左滑，受到向右的滑动摩擦力
$f_k=\mu_kmg$，做匀加速运动。若传送带足够长，且物块不受其他水平外力，当
$v_{\text{物}}=v_{\text{带}}$ 后，物块与带面相对静止，此后保持匀速，摩擦力为零。

\[
\text{加速阶段：}\quad
a=\mu_kg,\qquad
v_{\text{物}}=\mu_kgt,\qquad
x_{\text{物}}=\frac12\mu_kgt^2
\]
达到共速所需的时间和位移为
\[
t_1=\frac{v_{\text{带}}}{\mu_kg},
\qquad
x_1=\frac{v_{\text{带}}^2}{2\mu_kg}.
\]
若传送带有效长度小于 $x_1$，物块将在尚未共速时离开传送带。

\begin{center}
% TikZ 图：水平传送带模型
% Day 07：连接体、斜面与摩擦综合
\begin{tikzpicture}[scale=0.9,>=Latex]
  % 传送带
  \draw[thick] (0,0) circle (0.45);
  \draw[thick] (5,0) circle (0.45);
  \draw[thick] (0,0.45) -- (5,0.45);
  \draw[thick] (0,-0.45) -- (5,-0.45);
  \draw[->,thick,blue] (3.05,0.72) -- (4.25,0.72) node[right] {$v_{\text{带}}$};

  % 物块放在上表面
  \draw[thick] (1.2,0.45) rectangle (2.4,1.35);
  \node at (1.8,0.9) {$m$};
  \node[above] at (1.8,1.35) {$v_{\text{物}}=0$};

  % 传送带对物块的摩擦力
  \draw[->,thick,red] (2.4,0.9) -- (3.55,0.9) node[above] {$f_k$};
  \node[font=\footnotesize] at (2.55,-0.85) {刚放上时，物块相对传送带向左滑，$f_k$ 向右};
\end{tikzpicture}


\figcaption{传送带：物块刚放上时受向右摩擦力加速}\label{fig:conveyor-belt}
\end{center}

\subsection*{模型 4：倾斜传送带}

设长度为 $L$ 的传送带\textbf{沿斜面向下}以速度 $v_0$ 匀速运行，物块从斜面顶端由静止轻放；取\textbf{沿斜面向下为正}。这个模型必须分阶段：
\begin{enumerate}
  \item \textbf{未共速：} 物块相对带面向上滑，滑动摩擦力沿斜面向下，且重力沿斜面的分量也向下：
\[
mg\sin\theta+f_k=ma_1,\qquad
f_k=\mu_kmg\cos\theta,
\]
故
\[
\boxed{a_1=g(\sin\theta+\mu_k\cos\theta)},\qquad
t_1=\frac{v_0}{a_1},\qquad
s_1=\frac{v_0^2}{2a_1}.
\]
只有当 $L\ge s_1$ 时，物块才能在离开传送带前达到共速；若
$L<s_1$，应直接按第一阶段运动到斜面底端。
  \item \textbf{刚共速：} 到 $t_1$ 时，物块速度恰为 $v_0$，但这只是相对速度暂时为零，\textbf{不能立刻断定摩擦力为零}。
  \item \textbf{共速后：} 若要继续与传送带共速，沿斜面合力应为零，需要静摩擦力向上，大小为 $f_s=mg\sin\theta$。因此必须满足
\[
mg\sin\theta\leq\mu_smg\cos\theta
\quad\Longleftrightarrow\quad
\boxed{\tan\theta\leq\mu_s}.
\]
若满足，物块以后以 $v_0$ 匀速下滑，摩擦力是\textbf{向上}的静摩擦力；若 $\tan\theta>\mu_s$，静摩擦力不足以维持共速，物块又会相对带面向下滑，此后滑动摩擦力向上：
\[
mg\sin\theta-\mu_kmg\cos\theta=ma_2,\qquad
\boxed{a_2=g(\sin\theta-\mu_k\cos\theta)}.
\]
通常有 $\mu_s\ge\mu_k$，此时由 $\tan\theta>\mu_s$ 可知 $a_2>0$；若题目把两个摩擦因数作为独立数据给出，应由上式单独判断 $a_2$ 的符号。
\end{enumerate}

\textbf{符号提醒：} 严格地说，能否维持共速应比较 $\tan\theta$ 与\textbf{静摩擦因数} $\mu_s$；一旦再次相对滑动，才改用动摩擦因数 $\mu_k$。高中许多题目只给一个摩擦因数 $\mu$，并近似取 $\mu_s=\mu_k=\mu$，这时才可直接比较 $\tan\theta$ 与 $\mu$。

\begin{center}
% TikZ 图：向下运行的倾斜传送带初始阶段
% Day 07：连接体、斜面与摩擦综合
\begin{tikzpicture}[scale=0.98,>=Latex,line cap=round,line join=round]
  \def\ang{22}

  % 带面及其运动方向。
  \draw[thick] (0,0) -- (5.70,2.31);
  \draw[thick] (0,-0.28) -- (5.70,2.03);
  \foreach \x in {0.60,1.45,2.30,3.15,4.00}
    \draw[->,blue,thick] (\x,{0.405*\x-0.14}) -- ++(202:0.48);

  \draw[->,ultra thick,blue] (5.20,2.58) -- ++(202:1.15)
    node[midway,above=4pt] {$v_{\text{带}}$};

  % 物块下表面与带面贴合。
  \coordinate (P) at (2.36,0.96);
  \coordinate (C) at ($(P)+(112:0.36)$);
  \begin{scope}[shift={(P)},rotate=22]
    \draw[fill=white,thick] (-0.55,0) rectangle (0.55,0.72);
    \node at (0.27,0.48) {$m$};
  \end{scope}

  % 初始阶段的三个真实力，均从物块中心出发。
  \draw[->,very thick,red] (C) -- ++(202:1.12)
    node[below left=1pt] {$f_k$};
  \draw[->,very thick,blue] (C) -- ++(0,-1.28)
    node[below right=1pt] {$mg$};
  \draw[->,very thick,blue] (C) -- ++(112:1.02)
    node[above left=1pt] {$N$};

  % 倾角与阶段说明。
  \draw[thin] (0,0) -- (1.25,0);
  \draw (0.72,0) arc[start angle=0,end angle=22,radius=0.72];
  \node at (0.98,0.17) {$\theta$};
  \node[font=\footnotesize,align=center] at (3.15,-0.72)
    {初始时物块相对带面向上滑\\
     滑动摩擦力沿斜面向下};
\end{tikzpicture}


\figcaption{倾斜传送带初始阶段：带面向下运动，物块由静止释放，摩擦力沿斜面向下}\label{fig:inclined-conveyor-belt}
\end{center}

\section*{六、例题精讲}

\begin{examplebox}{水平连接体——整体法与隔离法}

\textbf{题意分析：} $m_A = 4\ \mathrm{kg}$、$m_B = 2\ \mathrm{kg}$ 用轻绳连接放在粗糙水平面上，并已经在拉力作用下向右滑动，动摩擦因数 $\mu_k = 0.3$，$g = 10\ \mathrm{m/s^2}$。在 $A$ 上施加 $F = 30\ \mathrm{N}$ 的水平拉力。求：
\begin{enumerate}
  \item 整体的加速度。
  \item 绳中张力 $T$。
\end{enumerate}

\textbf{物理模型：} 水平面粗糙连接体。

\textbf{解题过程：}

\begin{enumerate}
  \item 整体法：

  总摩擦力 $f_{\text{总}} = \mu_k (m_A + m_B) g = 0.3 \times 6 \times 10 = 18\ \mathrm{N}$

  \[
  F - f_{\text{总}} = (m_A + m_B) a,\qquad 30 - 18 = 6a,\qquad a = 2\ \mathrm{m/s^2}
  \]

  \item 隔离 $B$（求张力 $T$）：

  $B$ 受拉力 $T$（向右）和摩擦力 $f_B = \mu_k m_B g = 0.3 \times 2 \times 10 = 6\ \mathrm{N}$（向左）。

  \[
  T - f_B = m_B a,\qquad T - 6 = 2 \times 2,\qquad T = 10\ \mathrm{N}
  \]

  验证（隔离 $A$）：$F - T - f_A = m_A a$，$30 - 10 - 12 = 8 = 4 \times 2$，正确。
\end{enumerate}

\textbf{方法总结：}
\begin{enumerate}
  \item 先用整体法求共同加速度：把连接体看成一个系统，内力（绳中张力）不写，只写外力和总摩擦力。
  \item 再用隔离法求内力：选受力较少的物体列方程，未知量会更少。
  \item 最后用另一个物体反代验证。若两个隔离方程给出的张力不一致，通常是摩擦力方向或某个质量代错。
\end{enumerate}

\end{examplebox}

\begin{examplebox}{斜面摩擦临界}

\textbf{题意分析：} 质量为 $m$ 的物体放在倾角 $\theta = 30^\circ$ 的斜面上，$g = 10\ \mathrm{m/s^2}$。问：
\begin{enumerate}
  \item 求物体保持静止所需静摩擦因数 $\mu_s$ 的最小值（临界值）。
  \item 若 $\mu_s = 0.4$，物体会下滑吗？
  \item 若物体匀速上滑，且 $m = 2\ \mathrm{kg}$、$\mu_k = 0.4$，求所需的沿斜面向上的力 $F$。
\end{enumerate}

\textbf{物理模型：} 斜面摩擦临界 + 有外力时的动力学。

\textbf{解题过程：}

\begin{enumerate}
  \item 极限静止时 $\tan\theta = \mu_{s,\min}$：
  \[
  \mu_{s,\min} = \tan 30^\circ = \frac{\sqrt{3}}{3} \approx 0.577
  \]

  \item $\mu_s = 0.4 < 0.577$，最大静摩擦力不够，物体会下滑。

  \item 匀速上滑时，$a = 0$，滑动摩擦力沿斜面向下：
  \[
  F = mg\sin\theta + \mu_k mg\cos\theta
  \]
  \[
  F = mg(0.5 + 0.4 \times 0.866)
  \]
  \[
  F = 2 \times 10 \times 0.846 \approx 16.9\ \mathrm{N}
  \]
\end{enumerate}

\textbf{方法总结：}
\begin{enumerate}
  \item 静止临界先写 $mg\sin\theta=f_{s\max}$、$N=mg\cos\theta$，再得到 $\tan\theta=\mu_s$，不要把它当成无条件公式。
  \item 判断能否静止看 $\tan\theta$ 与 $\mu_s$；一旦已经滑动，计算摩擦力要改用 $\mu_k N$。
  \item 每换一种运动趋势，都要重新判断摩擦力方向。匀速上滑时摩擦力沿斜面向下；若物体有下滑趋势，静摩擦力沿斜面向上。
\end{enumerate}

\end{examplebox}

\begin{examplebox}{传送带模型——速度图像求相对滑动}

\textbf{题意分析：} 水平传送带以 $v_{\text{带}} = 4\ \mathrm{m/s}$ 匀速运行，将一质量为 $m = 2\ \mathrm{kg}$ 的物块轻放在传送带左端，$\mu_k = 0.4$，$g = 10\ \mathrm{m/s^2}$，传送带足够长。求：
\begin{enumerate}
  \item 物块达到传送带速度所需时间和位移。
  \item 物块在传送带上留下的摩擦痕迹长度。
\end{enumerate}

\textbf{物理模型：} 传送带——物块先匀加速后匀速。

\begin{center}
% Day 07：水平传送带例题图
\begin{tikzpicture}[scale=0.78,>=Latex]
  \draw[thick] (0,0) circle (0.36);
  \draw[thick] (5.2,0) circle (0.36);
  \draw[thick] (0,0.36) -- (5.2,0.36);
  \draw[thick] (0,-0.36) -- (5.2,-0.36);
  \draw[->,ultra thick,blue] (3.38,0.66) -- (4.55,0.66)
    node[right] {$v_{\text{带}}=4\ \mathrm{m/s}$};
  \draw[fill=white,thick] (0.78,0.36) rectangle (1.82,1.03);
  \node at (1.30,0.70) {$m$};
  \draw[->,thick,red] (1.82,0.70) -- (2.84,0.70) node[above] {$f_k$};
  \draw[->,thick,orange] (1.30,1.38) -- (0.35,1.38)
    node[midway,above] {$v_{\text{相对}}$};
  \node[font=\footnotesize] at (2.60,-0.72)
    {刚放上时：$v_{\text{物}}=0$};
\end{tikzpicture}


\figcaption{例题传送带：物块刚放上时相对带面向左滑}\label{fig:example-conveyor-scratch}
\end{center}

\textbf{解题过程：}

\begin{enumerate}
  \item 加速阶段：
  \[
  a = \mu_k g = 0.4 \times 10 = 4\ \mathrm{m/s^2}
  \]
  \[
  v_{\text{物}} = a t,\qquad
  t_1 = \frac{v_{\text{带}}}{a} = \frac{4}{4} = 1\ \mathrm{s}
  \]
  \[
  x_{\text{物}} = \frac12 a t_1^2 = \frac12 \times 4 \times 1^2 = 2\ \mathrm{m}
  \]

  \item 把物块视为质点时，痕迹长度等于物块与传送带的相对位移：
  \[
  x_{\text{带}} = v_{\text{带}} t_1 = 4 \times 1 = 4\ \mathrm{m}
  \]
  \[
  \Delta x = x_{\text{带}} - x_{\text{物}} = 4 - 2 = 2\ \mathrm{m}
  \]

  也可直接画 $v$-$t$ 图：带速为 $4\ \mathrm{m/s}$ 的水平线，物块速度从 $0$ 匀加速到 $4\ \mathrm{m/s}$。两图线间的面积等于相对位移：
  \[
  s_{\text{相对}}=\frac12\times 1\times4=2\ \mathrm{m}.
  \]
  \begin{center}
  % Day 07：传送带相对位移的 v-t 图
\begin{tikzpicture}[scale=0.76,>=Latex]
  \draw[->] (0,0) -- (5.45,0) node[right] {$t$};
  \draw[->] (0,0) -- (0,3.40) node[above] {$v$};
  \draw[very thick,blue] (0,2.55) -- (4.85,2.55)
    node[above right] {$v_{\text{带}}=4$};
  \draw[very thick,red] (0,0) -- (2.25,2.55)
    node[above left] {$v_{\text{物}}=4t$};
  \fill[orange!25] (0,0) -- (0,2.55) -- (2.25,2.55) -- cycle;
  \draw[orange!75!black,thick] (0,0) -- (0,2.55) -- (2.25,2.55);
  \draw[dashed] (2.25,0) -- (2.25,2.55);
  \node[below] at (2.25,0) {$1\ \mathrm{s}$};
  \node[font=\footnotesize,blue!70!black] at (3.55,2.22)
    {$t\ge1\ \mathrm{s}$ 时两者共速};
  \node[orange!85!black,align=center,font=\footnotesize] at (0.82,1.55)
    {$s_{\text{相对}}$\\$=\frac12\times1\times4$};
\end{tikzpicture}


  \figcaption{$v$-$t$ 图中两条速度线之间的面积等于相对位移}\label{fig:conveyor-vt-scratch}
  \end{center}
\end{enumerate}

\textbf{方法总结：}
\begin{enumerate}[itemsep=0.65em,topsep=0.4em,parsep=0pt]
  \item 传送带问题的核心是判断摩擦力方向（相对运动或相对运动趋势方向）。
  \item 物块刚放上时相对传送带向后滑，摩擦力向前——这是动力。
  \item 两条 $v$-$t$ 图线之间的面积也可求相对位移。\textbf{只有把物块视为质点、接触区域可忽略时}，带面上被单个接触点扫过的痕迹长度才等于相对位移；若物块底面有有限长度，题目所说“痕迹长度”还要结合接触区域的初末位置判断。
\end{enumerate}

\end{examplebox}

\section*{七、变式训练}

\begin{enumerate}[itemsep=0.65em,topsep=0.4em,parsep=0pt]
  \item $m_A = 5\ \mathrm{kg}$、$m_B = 3\ \mathrm{kg}$ 用轻绳连接放在光滑水平面上，在 $A$ 上施加 $F = 24\ \mathrm{N}$。求加速度和绳中张力。

  \item 质量为 $m = 3\ \mathrm{kg}$ 的物体放在倾角 $\theta = 37^\circ$（$\sin\theta = 0.6$，$\cos\theta = 0.8$）的斜面上，$\mu_s = 0.5$，若下滑则取 $\mu_k = 0.5$，$g = 10\ \mathrm{m/s^2}$。
  \begin{enumerate}[itemsep=0.35em,topsep=0.25em,parsep=0pt]
    \item 物体是否下滑？
    \item 若下滑，加速度多大？若静止，摩擦力多大？
  \end{enumerate}

  \item 水平传送带长 $L = 10\ \mathrm{m}$，以
  $v_{\text{带}} = 6\ \mathrm{m/s}$ 运行，$\mu_k = 0.3$，
  $g = 10\ \mathrm{m/s^2}$。物块从传送带左端轻放。判断物块到达右端前是否已与传送带共速，并求到达右端所需时间及相对滑动距离。
  \begin{center}
  \begin{tikzpicture}[scale=0.82,>=Stealth]
  \draw[thick] (0,0) rectangle (5.2,0.42);
  \foreach \x in {0.35,1.1,1.85,2.6,3.35,4.1,4.85}
    \draw[->,blue] (\x,0.21) -- ++(0.45,0);
  \draw[fill=gray!12] (0.35,0.42) rectangle (1.05,0.92);
  \draw[->,red,thick] (1.05,0.92) -- (1.8,0.92) node[midway,above] {$f_k$};
  \draw[<->] (0,-0.35) -- (5.2,-0.35) node[midway,below] {$L$};
  \node[font=\small] at (3.65,0.88) {带速 $v_{\text{带}}$};
  \node[below] at (0.35,0) {左端轻放};
\end{tikzpicture}

  \end{center}
\end{enumerate}

\begin{learningbox}{变式训练参考结果}
\begin{enumerate}[itemsep=0.45em,topsep=0.3em,parsep=0pt]
  \item 光滑水平连接体：
  \[
  a=\frac{F}{m_A+m_B}=3.0\ \mathrm{m/s^2},
  \qquad
  T=m_Ba=9.0\ \mathrm{N}.
  \]

  \item 斜面物块：$\tan37^\circ=0.75>\mu_s=0.5$，物块下滑。滑动摩擦力
  $f_k=\mu_kmg\cos37^\circ=12\ \mathrm{N}$，加速度
  \[
  a=g(\sin37^\circ-\mu_k\cos37^\circ)
   =2.0\ \mathrm{m/s^2}.
  \]

  \item 水平传送带：$a=\mu_kg=3.0\ \mathrm{m/s^2}$，
  $t_1=v_{\text{带}}/a=2.0\ \mathrm{s}$，
  $x_1=\frac12at_1^2=6.0\ \mathrm{m}<L$，所以到达右端前已经共速。总时间
  \[
  t=t_1+\frac{L-x_1}{v_{\text{带}}}
   =\frac83\ \mathrm{s},
  \]
  相对滑动距离为两条 $v$-$t$ 图线在 $0$ 到 $t_1$ 间的面积：
  $s_{\text{相对}}=\frac12v_{\text{带}}t_1=6.0\ \mathrm{m}$。
\end{enumerate}
\end{learningbox}

\section*{八、竞赛入门补充}

\begin{competitionbox}
\textbf{运动方向不同且含摩擦的连接体}

如果连接体中各物体的加速度方向不同（例如跨滑轮连接、斜面上下滑连接），不能直接把所有质量相加写成一个方向上的整体方程，通常要分别隔离。

\textbf{例题：} 质量 $m_A$ 和 $m_B$ 通过跨过光滑定滑轮的轻绳连接，$A$ 在粗糙水平面上，动摩擦因数为 $\mu_k$。系统已经运动，$B$ 竖直向下。求加速度和绳中张力。

\begin{center}
% TikZ 图：定滑轮连接体动力学
% Day 07：连接体、斜面与摩擦综合
\begin{tikzpicture}[scale=0.92,>=Latex]
  \path[use as bounding box] (-0.4,-2.95) rectangle (4.85,1.20);
  \coordinate (O) at (3.0,0);
  \def\r{0.34}

  % Table and wall corner; pulley center is fixed at the corner.
  \draw[thick] (-0.25,0) -- (O) -- (3.0,-2.35);
  \foreach \x in {0.0,0.35,...,2.65} {
    \draw[gray!60] (\x,0) -- ++(-0.16,-0.13);
  }
  \foreach \y in {-0.25,-0.60,...,-2.10} {
    \draw[gray!60] (3.0,\y) -- ++(0.13,0.17);
  }

  % Blocks.
  \draw[fill=white,thick] (0.75,0) rectangle (1.55,0.58);
  \node at (1.15,0.30) {$m_A$};
  \draw[fill=white,thick] (3.06,-1.65) rectangle (3.68,-0.84);
  \node at (3.37,-1.25) {$m_B$};

  % Pulley and one continuous light string.
  \draw[fill=white,thick] (O) circle (\r);
  \fill (O) circle (0.035);
  \draw[thick] (1.55,\r) -- ($(O)+(0,\r)$)
    arc[start angle=90,end angle=0,radius=\r]
    -- ($(O)+(\r,-0.84)$);

  % Forces and accelerations.
  \draw[->,thick,red] (1.55,0.29) -- (2.22,0.29) node[midway,below] {$T$};
  \draw[->,thick,blue] (0.95,0.83) -- (1.78,0.83) node[right] {$a$};
  \draw[->,thick,red] (3.37,-0.84) -- (3.37,-0.22) node[midway,left] {$T$};
  \draw[->,thick,blue] (3.37,-1.65) -- (3.37,-2.30) node[midway,left] {$m_Bg$};
  \draw[->,thick,blue] (4.08,-1.18) -- (4.08,-1.88) node[midway,right] {$a$};

  \node[font=\footnotesize,align=center] at (1.85,-2.62)
    {轻绳不伸长：$A$ 向右与 $B$ 向下的加速度大小相同};
\end{tikzpicture}


\figcaption{滑轮连接体：两物体加速度大小相同、方向不同}\label{fig:pulley-connected-dynamics}
\end{center}

\textbf{解：} 轻绳不伸长，所以两物体加速度大小相同，但方向不同。分别隔离：
\begin{itemize}
  \item 对 $A$（向右为正）：$T-\mu_km_Ag=m_Aa$
  \item 对 $B$（竖直方向）：$m_B g - T = m_B a$
\end{itemize}

联立得：
\[
a=\frac{m_B-\mu_km_A}{m_A+m_B}g,
\qquad
T=\frac{(1+\mu_k)m_Am_B}{m_A+m_B}g.
\]

\textbf{条件检查：} 上式按 $B$ 向下取正。题目已说明系统按该方向运动；若只说“由静止释放”，还应先用最大静摩擦力判断能否启动。轻绳、光滑滑轮保证绳中各处张力大小相等。
\end{competitionbox}

\begin{competitionbox}
\textbf{速度约束：先写整根绳的长度关系}

有些连接体题中，绳子不是沿物体运动方向拉动物体。下图中，轻绳一端连接物块，穿过固定光滑小环后，另一端竖直向下运动。设斜绳段长为 $\ell$、竖直绳段长为 $y$，忽略小环处的定长部分，则
\[
L=\ell+y+\text{定长部分}.
\]

\begin{center}
% TikZ 图：斜拉绳的速度约束
% Day 07：连接体、斜面与摩擦综合
\begin{tikzpicture}[scale=0.95,>=Latex]
  \path[use as bounding box] (-0.75,-1.35) rectangle (5.15,2.25);

  \coordinate (P) at (0.15,0);
  \coordinate (R) at (4.05,1.55);
  \coordinate (Q) at (1.20,0.42);
  \coordinate (H) at (1.35,0);

  % Ground / waterline and fixed point.
  \draw[thick] (-0.55,-0.36) -- (2.35,-0.36);
  \foreach \x in {-0.40,-0.05,...,2.20} {
    \draw[gray!60] (\x,-0.36) -- ++(-0.16,-0.12);
  }
  \draw[thick] (4.05,0.55) -- (4.05,2.05);
  \foreach \y in {0.65,0.98,...,1.95} {
    \draw[gray!60] (4.05,\y) -- ++(0.15,0.10);
  }
  \draw[fill=white,thick] (R) circle (0.12);
  \node[above right=2pt] at (R) {固定点};

  % Block and rope.
  \draw[fill=white,thick] (-0.23,-0.02) rectangle (0.53,0.42);
  \node at (0.15,0.20) {$m$};
  \draw[thick] (P) -- (R);

  % Horizontal velocity and reference line.
  \draw[dashed,gray] (P) -- (1.75,0);
  \draw[->,thick,blue] (0.15,-0.63) -- (1.45,-0.63)
    node[midway,below] {$v$};

  % Projection along the rope.
  \draw[->,very thick,orange] (P) -- (Q)
    node[pos=0.58,above left=2pt] {$v_{\parallel}=v\cos\theta$};
  \draw[->,thick,red] (3.05,1.15) -- (3.70,1.41)
    node[midway,below right=1pt] {$u$};

  % Angle marker.
  \draw (0.56,0) arc[start angle=0,end angle=21.7,radius=0.41];
  \node at (0.86,0.16) {$\theta$};

  % Short explanatory formula.
  \node[align=center] at (2.35,-1.02)
    {绳长变化只看沿绳方向的分量：\quad $u=v\cos\theta$};
\end{tikzpicture}


\figcaption{斜拉绳的速度约束：绳长变化只由沿绳方向的速度分量决定}\label{fig:rope-velocity-projection}
\end{center}

在很短时间 $\Delta t$ 内，物块向右移动，斜绳段缩短量近似为
$v\cos\theta\,\Delta t$；为保持整根绳长不变，竖直绳段必须增加同样的长度。若自由端向下速度为 $u$，则
\[
u\Delta t=v\cos\theta\,\Delta t,
\qquad
\boxed{u=v\cos\theta}.
\]

用瞬时变化率也可写成
\[
-\frac{\mathrm d\ell}{\mathrm dt}
=\frac{\mathrm dy}{\mathrm dt}
=u=v\cos\theta.
\]
\end{competitionbox}

\section*{九、易错点}

\begin{mistakebox}
\begin{enumerate}
  \item \textbf{整体法漏了摩擦力：} 先把系统外力逐个圈出。粗糙表面上的外摩擦力要写进整体方程，绳张力这类内力不写。

  \item \textbf{隔离法选错对象：} 先选择未知力少、受力简单的物体隔离。水平连接体中常先隔离后面的物体求张力。

  \item \textbf{传送带问题中摩擦力方向判断错：} 先判断物块相对传送带的滑动或趋势。刚放上时常相对向后滑，摩擦力向前；共速后再重新判断。

  \item \textbf{把 $\tan\theta$ 与 $\mu_s$ 的比较当成通用口诀：} 这条比较只适用于物体静放在斜面上、沿斜面方向没有其他外力的情形。先假设物体将要下滑，把静摩擦力画成沿斜面向上；再由临界平衡写 $mg\sin\theta=f_{\max}=\mu_smg\cos\theta$，才得到 $\tan\theta=\mu_s$。有拉力、推力或其他约束时，应重新列沿斜面的平衡方程。

  \item \textbf{匀速上滑时把摩擦力方向画反：} 匀速上滑时物体相对斜面向上运动，摩擦力沿斜面向下（阻力），不是向上。

  \item \textbf{斜拉绳时直接写速度相等：} 如下图，轻绳穿过固定光滑小环，斜绳段长为 $\ell$、竖直段长为 $y$。由
  $L=\ell+y+\text{定长部分}$ 可知
  $u=-\mathrm d\ell/\mathrm dt=v\cos\theta$，而不是
  $u=v$。改变装置或增加绳段后，必须重新写整根绳的长度关系。
\end{enumerate}

\begin{center}
% Day 07：斜拉绳速度约束易错辨析
\begin{tikzpicture}[scale=0.73,>=Latex,line cap=round,line join=round]
  \path[use as bounding box] (-0.55,-0.92) rectangle (7.15,2.18);
  \coordinate (P) at (0.38,0.18);
  \coordinate (R) at (3.30,1.34);
  \coordinate (C) at (3.30,-0.42);

  % 合法的绳长约束装置：绳穿过固定小环。
  \draw[thick] (-0.42,-0.28) -- (1.28,-0.28);
  \draw[fill=white,thick] (-0.38,-0.02) rectangle (0.38,0.38);
  \node at (0,0.18) {$m$};
  \draw[thick] (3.62,-0.70) -- (3.62,1.82);
  \draw[thick] (R) -- (3.62,1.34);
  \draw[very thick] (P) -- (R) -- (C);
  \draw[fill=white,thick] (R) circle (0.11);
  \fill (C) circle (2pt);
  \node[above right=1pt] at (R) {固定小环};

  \draw[->,thick,blue] (-0.02,-0.60) -- (1.05,-0.60)
    node[midway,above=2pt] {$v$};
  \draw[->,thick,blue] (3.92,0.10) -- (3.92,-0.62)
    node[midway,right] {$u$};
  \draw[->,thick,orange] (P) -- ($(P)!0.34!(R)$);
  \draw[dashed,gray] (P) -- ++(0.95,0);
  \draw (0.72,0.18) arc[start angle=0,end angle=21.7,radius=0.34];
  \node at (0.94,0.31) {$\theta$};

  % 错误与正确关系分开放置。
  \node[font=\footnotesize,red!80!black] at (5.70,1.32)
    {错误：$u=v$};
  \draw[red!80!black,very thick] (5.08,0.98) -- (6.32,1.66);
  \draw[red!80!black,very thick] (5.08,1.66) -- (6.32,0.98);
  \node[font=\footnotesize,green!50!black] at (5.70,0.28)
    {正确：$u=v\cos\theta$};
  \node[font=\footnotesize] at (1.95,-0.72)
    {$L=\ell+y+\text{定长部分}$};
\end{tikzpicture}


\figcaption{斜拉绳易错辨析：先看绳长变化，再比较速度分量}\label{fig:rope-velocity-mistake}
\end{center}
\end{mistakebox}

\section*{十、今日小结}

\begin{enumerate}[nosep]
  \item 整体法：求共同加速度时用，只考虑系统外力。
  \item 隔离法：求内力时用，先整体求加速度再隔离。
  \item 斜面临界：无其他沿斜面外力时，$\tan\theta = \mu_s$ 表示物块处于极限静止状态。
  \item 传送带：先检查物块能否在离开带面前共速；水平带上共速后可无摩擦，斜带上共速后仍要检查静摩擦力能否平衡重力沿斜面的分量。
  \item 摩擦方向检查：先判断相对运动或相对运动趋势；摩擦力与相对运动（趋势）方向相反，不是简单与运动方向相反。
  \item 斜拉绳的速度约束要先写整根绳长关系，再看沿绳方向的速度分量；图示小环模型中 $u=v\cos\theta$。
\end{enumerate}

\textbf{一分钟复盘：} 看到连接体题，先判断“是否共同加速度、是否要求内力、整根绳由哪些可变绳段组成”。这三问回答完，再决定整体法、隔离法或速度投影。
